\section{Related Work}
\label{sec:related}

\para{Fine-tuning degrades safety.} A now-robust line of work shows that minimal
fine-tuning --- even ten examples, even benign data --- can strip the safety alignment of
an aligned model \citep{qi2023finetuning}. \citet{qi2024shallow} explain this as
\emph{shallow} safety alignment: alignment mostly adapts the distribution of the first few
output tokens, so fine-tuning, jailbreaks, and prefilling all break it easily. This
motivates our choice to measure drift over full generated sequences rather than first
tokens, and frames \emph{why} robustness, where it exists, may be narrow.

\para{Emergent misalignment and its controls.} \citet{betley2025em} show that narrow
fine-tuning on insecure code induces \emph{broad} misalignment on unrelated prompts,
strongest in large models (GPT-4o, Qwen2.5-Coder-32B) and notably \emph{inconsistent}
($\sim$20\% of answers). Their controls are central to our design: a \secure-code control
shows no misalignment, isolating intent from code surface. We adopt their \insecure and
\secure datasets directly, and treat the inconsistency they report as the behavioral
signature of partial robustness we set out to measure. \citet{soligo2026em} show a single
rank-1 LoRA direction can mediate emergent misalignment, justifying our use of low-rank
fine-tuning, while \citet{mushtaq2025narrow} demonstrate that narrow unlearning can also
spill over into broad misalignment.

\para{The geometry of alignment in weight space.} \citet{peng2024landscape} discover a
parameter-space ``safety basin'': random weight perturbations within a local neighborhood
preserve safety, then safety drops sharply outside it, and they propose the VISAGE metric.
\citet{yang2025asft} find that perturbations \emph{orthogonal} to the alignment direction
rapidly compromise safety, framing the basin as narrow. \citet{wei2024brittleness} localize
safety to a sparse ($\sim$3\%) set of parameters. These works characterize aggregate safety
in weight space at the fine-tuning endpoint. Unlike them, we work in \emph{behavior} space,
measuring per-prompt output trajectories across checkpoints, and we vary the
\emph{semantic distance} of the probe rather than the weight-perturbation direction.

\para{Is orthogonality stable?} \citet{springer2026geometry} argue that the reassurance
``fine-tuning updates are orthogonal to safety directions, so safety is preserved'' is
false: orthogonality is structurally unstable under gradient descent and collapses, with
alignment concentrated in low-dimensional subspaces. \citet{hsiung2025similarity} find that
safety collapse correlates with the representational similarity between alignment and
fine-tuning data --- direct support for a distance-dependent view. Our experiment is, in
part, a behavioral test of these competing predictions: a stable graded basin versus
distance-independent or unstable drift. Our results align with the instability view.

\para{Training as a trajectory.} \citet{danovski2024dynamical} treat training as a dynamical
system and measure stability and chaos along the trajectory, directly motivating our
``trajectory attractor'' framing and reversal-rate metric. \citet{rahimi2026linearization}
show that with a stay-close-to-pretrained bias, fine-tuning dynamics become equivalent to
NTK learning, implying an implicit attraction toward the base model.
\citet{ilharco2022taskarithmetic} establish task vectors (the weight difference
$\vtheta_{\text{ft}}-\vtheta_{\text{pre}}$) as meaningful directions, the vocabulary for
weight-space divergence. \citet{zhang2024unforgettable} show that whether forgetting
generalizes is highly task-dependent --- the closest prior demonstration of \emph{latent}
robustness to fine-tuning and its task-dependence. We combine the trajectory view with the
emergent-misalignment intervention to build a \emph{predictive} per-prompt divergence
diagnostic, which to our knowledge is new.
